% !TEX root = chapter-standalone.tex


\section{Optimality}
\label{sec:dp-optimality}

\todotext{2026-07, J: add intro material here}

\subsection{Chains and Antichains}
\label{sec:chains-antichains}


Here are two further special types of subsets of a poset: \SY{chains} and \SY{antichains}.
Their definitions are dual.

\begin{definition}[Chain in a poset]
    \label{def:chain}
    Given a poset~$\posAdefinition$, a \maindef{chain} is a subset~$\subA \setsubseteq \posAset$ such that any two elements of~$\subA$ are comparable:
    \begin{equation}\label{eq:chain}
        \prfperiod{
            \posela,\poselb \setin \subA
        }{
            (\posela \posAleq \poselb) \boolor (\poselb \posAleq \posela)
        }
    \end{equation}
\end{definition}


\begin{definition}[Antichain in a poset]
    \label{def:antichain}
    An \maindef{antichain} is a subset~$\subA$ of a \SY{poset} where no two distinct elements are comparable:
    %
    \begin{equation}\label{eq:antichain}
        \prfperiod{
            \posela, \poselb \setin \subA
        }{
            \posela \posAleq \poselb
        }{
            \posela = \poselb
        }
    \end{equation}
    %
\end{definition}
\begin{remark}
    Note that the empty set~$\Emptyset$ is both a \SY{chain} and an \SY{antichain}.
\end{remark}

We denote the set of \SY{antichains} of a poset~\posA by~$\antichains\posA$.

\begin{example}[Chains and antichains in a power poset]
    Consider the \SY{poset} in \cref{fig:powersetcat3}.
    Examples of \SY{chains} are
    \begin{equation}\label{eq:example-chain}
        \makeset{\Emptyset,\makeset{\posela}}
        \qqand
        \makeset{\Emptyset,\makeset{\poselb},\makeset{\poselb,\poselc},\makeset{\posela,\poselb,\poselc}},
    \end{equation}
    depicted in \cref{fig:power_chains_a} and \cref{fig:power_chains_b}, respectively.

    Examples of \SY{antichains} are
    \begin{equation}\label{eq:examples-antichain}
        \makeset{\makeset{\posela},\makeset{\poselb}}
        \qqand
        \makeset{\makeset{\posela,\poselb}, \makeset{\posela,\poselc}, \makeset{\poselb,\poselc}},
    \end{equation}
    depicted in \cref{fig:power_antichains_a} and \cref{fig:power_antichains_b}, respectively.
\end{example}

\begin{marginfigure}
    \subfloat[\label{fig:power_chains_a}
        A \SY{chain}.
    ]{
        \includesag{40_dpcatfig_power_chains}
    }

    \subfloat[\label{fig:power_chains_b}
        A \SY{chain}.
    ]{
        \includesag{40_dpcatfig_power_chains_bis}
    }

    \subfloat[\label{fig:power_antichains_a}
        An \SY{antichain}.
    ]{
        \includesag{40_dpcatfig_power_antichains}
    }

    \subfloat[\label{fig:power_antichains_b}
        An \SY{antichain}.
    ]{
        \includesag{40_dpcatfig_power_antichains_bis}
    }
    \caption{Examples of \SY{chains} (a-b) and \SY{antichains} (c-d) in the \SY{poset} $\powerset\makeset{\posela,\poselb,\poselc}$.}
\end{marginfigure}

\clearpage

\begin{marginfigure}
    \centering
    \includesag{70_antichain}
    \caption{Example of discrete \SY{antichains}.
    }
    \label{fig:antichain}
\end{marginfigure}
\begin{example}
    In the context of battery choices, consider the diagram in~\cref{fig:antichain}.
    The black markers represent an \SY{antichain} of choices
    \begin{equation}\label{eq:example-battery-antichain}
        \makeset{
            \tupp{\poscheap,\posheavy},
            \tupp{\posexpensive,\poslight}
        }.
    \end{equation}
    It is a set of pairs because they do not dominate each other: one is cheaper, but is heavier, and the other is more expensive, but lighter, making them incomparable.
    If a battery with the properties of the red marker existed (very expensive, between light and heavy), that would be an element that cannot be part of the \SY{antichain}, since it would be dominated by~$\tupp{\posexpensive,\poslight}$.
    %
    \begin{marginfigure}
        \centering
        \includesag{70_antichain_2}
        \caption{Example of continuous \SY{antichains}.}
        \label{fig:antichain_2}
    \end{marginfigure}
    %
    Similarly, we could think of a continuous law which relates battery cost and mass.
    Assume that cheap means~$\unit[10]{\standardcurrency}$, expensive means~$\unit[20]{\standardcurrency}$, light means~$\unit[250]{g}$, and heavy means~$\unit[500]{g}$.
    For instance, consider the \SY{antichain} given by~$\text{mass}=500-25\cdot \text{cost}$, with maximum possible cost~$\unit[20]{\standardcurrency}$ (\cref{fig:antichain_2}).
\end{example}

\vspace{5cm}

\begin{marginfigure}
    \centering
    \includesag{chain_divisor_poset}
    \caption{}
    \label{fig:chain_divisor_poset}
\end{marginfigure}

\begin{marginfigure}
    \centering
    \includesag{antichain_divisor_poset}
    \caption{}
    \label{fig:antichain_divisor_poset}
\end{marginfigure}

% We will now have a look at a couple more examples.
\begin{example}
    Consider the poset~$\posAdefinition$ where~$(\posAel \posAleq \posBel)$ if~$\posAel$ is a divisor of~$\posBel$ and~$\posAset=\makeset{1,5,10,11,13,15}$.
    A \SY{chain} of~\posA is~$\makeset{1,5,10}$ (\cref{fig:chain_divisor_poset}).
    An \SY{antichain} of~\posA is~$\makeset{10,11,13,15}$ (\cref{fig:antichain_divisor_poset}).
\end{example}

\devel{
\subsection{Measuring posets}

\todotext{decide if this subsection is a DP thing and if/where to put it}

We can define two measurements for a poset: the height and the width.
These measurements allow us to quantify the performance of several algorithms we will see in later parts of the book.
\begin{definition}[Width of a poset]
    \label{def:poset-width}
    The \maindef{width of a poset}, denoted~$\posetwidth(\posA)$, is the maximum cardinality of an \SY{antichain} in~\posA.
\end{definition}

\begin{definition}[Height of a poset]
    \label{def:poset-height}
    The \maindef{height of a poset}, denoted~$\posetheight(\posA)$, is the maximum cardinality of a \SY{chain} in~\posA.
\end{definition}

Note that an empty \SY{poset} has exactly one \SY{chain} and one \SY{antichain}: the empty set.
Therefore, the height and width are zero.

\begin{example}
    Consider the poset~\posA in \cref{fig:poset-height-width}.
    The longest \SY{antichains} of~\posA are~$\makeset{\sapple, \sburger, \scheese}$, $\makeset{\sapple, \schoco, \scheese}$, $\makeset{\sbretzel, \schoco, \scheese}$, $\makeset{\sapple, \schoco, \sgrapes}$, $\makeset{\sbretzel, \schoco, \sgrapes}$, and~$\makeset{\sbretzel, \sburger, \scheese}$.
    Therefore,
    \begin{equation}\label{eq:poset-height-width-1}
        \posetwidth(\posA)=3.
    \end{equation}
    The longest \SY{chain} in the \SY{poset} is~$\makeset{\scarrot,\schoco,\sburger,\sfondue}$, and therefore
    \begin{equation}\label{eq:poset-height-width-2}
        \posetheight(\posA)=4.
    \end{equation}
\end{example}

\begin{figure*}[h]
    \aligninner{
        \subfloat[
            \label{fig:poset-height-width}
        ]{
            \includesag{height_width_posets}
        }
        \hspace{3em}
        \subfloat[
            \label{fig:poset-height-width2}
            The longest \SY{chain} ]{
            \includesag{height_width_posets2}
        }
        \hspace{3em}
        \subfloat[
            \label{fig:poset-height-width3}
            One of the largest \SY{antichains} ]{
            \includesag{height_width_posets3}
        }
    }
    \caption{Example for height and width of a poset.
    }
\end{figure*}

\vfill
\begin{gradedexercise}[\exname{MeasurePowerPoset}]
    \label{ex:MeasurePowerPoset}

    Let~\setA be a finite set with~$n$ elements.
    Obtain an expression (without proof) for
    \begin{enumerate}
        \item $\posetwidth(\powerset \setA)$;
        \item $\posetheight(\powerset \setA)$.
    \end{enumerate}
\end{gradedexercise}

\solutionof{MeasurePowerPoset}

}

\subsection{Antichains and upper sets}



\begin{lemma}
    \label{lem:up-cl-inj-antichains}
    Let~\setA and~\setB be subsets of~\posA that are \SY{antichains}.
    Then
    \begin{equation}
        \prfperiod{
            \upit \setA = \upit \setB
        }{
            \setA = \setB
        }
        %\upit  \setA = \ \upit  \setB \quad \Imp \quad \setA = \setB.
    \end{equation}
\end{lemma}

\begin{proof}
    Fix an element~$\setAel \setin \setA$.
    From~$\upit \setA = \upit \setB$ we know that in particular~$\setA \setsubseteq \upit \setB$.
    This means that for our fixed~$\setAel \setin \setA$ there exists~$\setBel \setin \setB$ such that~$\setBel \posleq \setAel$.
    From~$\upit \setA = \upit \setB$ it also follows that~$\setB \setsubseteq \upit \setA$, so to the~$\setBel \setin \setB$ given above, there exists~$\setAel\styleelements{'} \setin \setA$ such that~$\setAel\styleelements{'} \posleq \setBel$.
    In total, we have~$\setAel\styleelements{'} \posleq \setBel \posleq \setAel$, and since~$\setA$ is an \SY{antichain}, we must have~$\setAel\styleelements{'} = \setAel$.
    This implies that~$\setAel\styleelements{'} = \setBel = \setAel$.
    In particular, we have~$\setAel \setin \setB$.

    The above shows that~$\setA \setsubseteq \setB$.
    To show~$\setB \setsubseteq \setA$, we can fix any~$\setBel \setin \setB$ and repeat the above argumentation, now with the roles of~\setA and~\setB exchanged.
\end{proof}



\begin{definition}[Downward closed set]
    \label{def:downward-closed-upperset}
    An \SY{upper set}~$\subA$ is \maindef{downward-closed} in a poset~\posA if
    \begin{equation}
        \subA = \upit \Min \subA.
    \end{equation}
    The set of downward-closed \SY{upper sets} of~\posA is denoted~$\setOfDCUppersets \posA$.

\end{definition}

\begin{definition}[Upward closed set]
    \label{def:upward-closed-lowerset}
    A \SY{lower set}~$\subA$ is \maindef{upward-closed} in a poset~\posA if
    \begin{equation}
        \subA = \downit \Max \subA.
    \end{equation}
    The set of upward-closed \SY{lower sets} of~\posA is denoted~$\setOfUCLowersets \posA$.
\end{definition}

